\documentclass[11pt]{article}
\usepackage{amsmath,amssymb,amsthm}
\newtheorem{lemma}{Lemma}
\newtheorem{proposition}{Proposition}
\theoremstyle{remark}\newtheorem*{remark}{Remark}
\title{Segmented inversion of the Gaussian correction and large-coefficient precision}
\date{Draft, October 8, 2026}
\begin{document}
\maketitle

\section*{Setting}
This uses PR~\#5's chirped-correlation lemma for the correlations and its parameter conventions. PR~\#5 is unmerged on \texttt{CrocSwap/integer-mult-bounds}.

Keep the notation of the one-dimensional resampling lemma:
$2\le s<t$, $\sigma=t/s$, $\theta=\sigma-1\le1/4$, $\alpha\ge2$ with $\alpha^2\theta\ge1$,
$q_j=\lfloor tj/s+1/2\rfloor$, $\beta_j=tj/s-q_j\in[-1/2,1/2)$, defined for all
integers $j$, so $q_{j+s}=q_j+t$ and $\beta_{j+s}=\beta_j$. The matrix $N=C\mathsf T\mathsf D$
acts on $s$-periodic sequences by
\[
 (Nu)_\ell=\sum_{j\in\mathbb Z}e^{-\pi\alpha^2X_{\ell,j}}u_j,\qquad
 X_{\ell,j}=(\sigma j-q_\ell)^2-\beta_j^2 .
\]

\section{Exact structure}

\begin{lemma}[Closed form]\label{lem:closed}
Put $n=q_j-q_\ell\in\mathbb Z$. Then $X_{\ell,j}=n(n+2\beta_j)$.
\end{lemma}
\begin{proof}
$\sigma j-q_\ell=(q_j+\beta_j)-q_\ell=n+\beta_j$, so
$X=(n+\beta_j)^2-\beta_j^2=n(n+2\beta_j)$.
\end{proof}

Since $\sigma=1+\theta$ with $\theta<1$, consecutive values satisfy
$q_{j+1}-q_j\in\{1,2\}$. Call $j$ a \emph{wrap} when the difference is $2$;
between wraps $\beta_{j+1}=\beta_j+\theta$. A \emph{segment} is a maximal run of
indices containing no wrap step. Each segment has $\lfloor1/\theta\rfloor$ or
$\lceil1/\theta\rceil$ indices, up to one boundary index, and wraps have density
$t/s-1=\theta$.

\begin{lemma}[Chirp--Toeplitz segments]\label{lem:segment}
Let $\ell,j$ lie in one segment and choose the real centre
$\ell^*=\ell-(\beta_\ell+1/2)/\theta$, which is the same for the whole segment.
With $u=\ell-\ell^*$, $v=j-\ell^*$ and $h=j-\ell$,
\[
 e^{-\pi\alpha^2X_{\ell,j}}
  =e^{\pi\alpha^2\theta u^2}\;\tau_h\;e^{-\pi\alpha^2\theta v^2},
 \qquad \tau_h=e^{-\pi\alpha^2(\sigma h^2-h)} .
\]
Thus the principal block $M_k$ of $N$ on segment $k$ equals $D_kT_{L_k}D_k^{-1}$, a similarity,
where $T_L$ is the $L\times L$ Toeplitz matrix with symbol $(\tau_h)$, which does not
depend on the segment, and $D_k,D'_k$ are diagonal.
\end{lemma}
\begin{proof}
Inside a segment $n=h$ and $\beta_j=\beta_\ell+\theta h=\theta v-1/2$. By
Lemma~\ref{lem:closed}, $X=h(h+2\theta v-1)$. Using $2hv=v^2-u^2+h^2$ gives
$X=(1+\theta)h^2-h+\theta v^2-\theta u^2$, which is the stated factorisation.
\end{proof}

\begin{lemma}[Cross-wrap entries]\label{lem:cross}
If $\ell,j$ lie in different segments then $|n|\ge|h|+1\ge2$ and
$X_{\ell,j}\ge|n|(|n|-1)\ge2$. In particular every cross-segment entry is at most
$e^{-2\pi\alpha^2}$. An entry exceeds $2^{-Q}$ only if $|n|(|n|-1)<Q/(\pi\alpha^2\log_2e)$.
\end{lemma}
\begin{proof}
Crossing $k\ge1$ wrap steps adds $k$ to $|q_j-q_\ell|$ beyond $|h|$. Lemma~\ref{lem:closed}
with $|\beta_j|\le1/2$ gives $X\ge|n|(|n|-1)$.
\end{proof}

The identities in Lemmas~\ref{lem:closed}--\ref{lem:cross} were checked in exact rational
arithmetic for $(s,t)=(113,128)$ and $(1009,1024)$, over three periods and $|j-\ell|\le30$:
no mismatches, and the minimum cross-segment exponent was $228/113$ and $2020/1009$.

\paragraph{Indexing.} Indices live in $\mathbb Z/s\mathbb Z$. Index $0$ has $\beta_0=0$, so it lies mid-segment, and one segment always crosses the period boundary; on tape it is stitched from the two line ends. Aliased terms $j+ms$ belong to $\Delta$, not to $M_k$. They matter only if one segment covers a whole period ($t-s=1$), which the parameters exclude.

\section{Applying $N^{-1}$}

Write $N=M+\Delta$ with $M=\bigoplus_kM_k$. Fix a target accuracy $2^{-Q}$ and drop
entries of $\Delta$ below $2^{-Q-\log_2 s-4}$. By Lemma~\ref{lem:cross}, what remains of
$\Delta$ sits in $w\times w$ corner blocks at each wrap, with
\[
 w\le1+\sqrt{Q/(4.53\,\alpha^2)} .
\]

\begin{proposition}[Segmented inverse]\label{prop:inverse}
Assume $\|N-I\|<1/2$, which holds when $\alpha^2\theta\ge1$. Then $N^{-1}u$ can be computed
to absolute error $2^{-Q}$, for $\|u\|\le1$, using
\begin{enumerate}
\item two applications of $M^{-1}$;
\item per wrap, $O(w^3)$ arithmetic operations at precision $P$ to form and factor the
capacitance block, and $O(w^2)$ to apply it;
\end{enumerate}
with working precision $P=Q+\lceil10.2\,\alpha^2/\theta\rceil+O(\log s)$.
Each application of $M^{-1}$ costs a constant number of chirped correlations of segment
length, hence $O(P^{1+\delta})$ per output. The setup computes the Gohberg--Semencul
vectors of $T_L$ and $T_{L+1}$ once per line, in $O(L^2)$ operations.
\end{proposition}

\begin{proof}[Proof sketch]
\emph{Segment blocks.} $M_k$ is a principal submatrix of $N=I+E$, so
$\|M_k-I\|\le\|E\|<1/2$ and $\|M_k^{-1}\|\le2$. Also
$\|T_L-I\|\le\sum_{h\ne0}\tau_h<2e^{-\pi\alpha^2\theta}<0.1$, so $T_L$ is well conditioned
and has a Gohberg--Semencul representation
$T_L^{-1}=x_0^{-1}\bigl(L(x)U(y)-L(z)U(w)\bigr)$ with triangular Toeplitz factors.
Applying $M_k^{-1}=D_k'^{-1}T_L^{-1}D_k^{-1}$ means: scale by
$e^{-\pi\alpha^2\theta u^2}\le1$, perform four triangular Toeplitz products (correlations,
evaluated as in the chirped-correlation lemma), then scale by
$e^{\pi\alpha^2\theta v^2}\le e^{\pi\alpha^2\theta(L+1)^2}$. Since $L\le1/\theta+1$, the
last factor is at most $2^{10.2\alpha^2/\theta}$ for $\theta\le1/4$ (and about $2^{4.53\alpha^2/\theta}$ as $\theta\to0$), and this is the precision
reserve in $P$. The exact result has norm at most $2\|u\|$.

\emph{Coupling.} Write $\Delta=UV^{\mathsf T}$, where $V$ selects the $2w$ columns
around each wrap and $U$ holds the corresponding columns of $\Delta$. Woodbury gives
\[
 N^{-1}=M^{-1}-M^{-1}U\,C^{-1}V^{\mathsf T}M^{-1},\qquad C=I+V^{\mathsf T}M^{-1}U .
\]
$C$ is cyclic block tridiagonal with $2w\times2w$ blocks, because $M^{-1}$ is block diagonal, each block of $U$ touches the two segments adjacent to its wrap, and indices are taken modulo $s$. This needs $2w<L$, which holds for $d\ge3$. Solve it by bordered (cyclic) block elimination, still $O(w^3)$ per wrap, with one forward and one backward sweep. Its entries need only
the corner entries of each $M_k^{-1}$. From the Gohberg--Semencul form, entry $(a,b)$ is a sum of $\min(a,b)+1$ products. For the bottom-right corner use persymmetry, $(T^{-1})_{L-1-a,L-1-b}=(T^{-1})_{b,a}$, so every corner entry costs $O(w)$, and $O(w^3)$ per wrap. Also
$\|C-I\|\le\|M^{-1}\|\,\|\Delta\|\le4we^{-2\pi\alpha^2}$, so block elimination without
pivoting is stable, at $O(w^3)$ per wrap.
$V^{\mathsf T}M^{-1}u$ reads entries of the first $M^{-1}u$. The vector $Uz$ is supported
near wraps, and the second $M^{-1}$ application handles it.
\end{proof}

\begin{remark}[Cost per output]
Wraps have density $\theta$. Per output the coupling therefore costs
$O(\theta w^3)$ operations at precision $P$.
\end{remark}

\section{Large coefficients}

Split the inputs into $b'$-bit digits instead of $b=\lceil\log_2n\rceil$-bit digits, with
$b'=\lceil b^{1+x}\rceil$ for a fixed $x>0$, and set $p=6b'$. Then $T\in[4n/b',8n/b')$,
the volume is $Tp=\Theta(n)$, and the address length is $\log_2T=b-(1+x)\log_2b+O(1)$.
Choose $d=\lfloor b^\epsilon\rfloor$ and $r$, $\ell$ from $T$ as before.
The capacity bound becomes $c_j<2^{3b'}$, and the final error bound becomes
\[
 2^{2b'+4}S\cdot28\cdot2^\gamma T^2L\cdot2^{-6b'}<448\,b'\,2^{-b'+\gamma}<\tfrac12
 \quad\text{whenever }\gamma\le b'/4 .
\]
With $\alpha^2=\lfloor b'/(8d)\rfloor$ we have $\gamma=2d\alpha^2\le b'/4$ and
$\alpha^2\theta\ge b'/(32d^2)\ge1$ once $2\epsilon<1+x$. All uses of
$r\ge2^{a_rp/d}$ in the layer and transform proofs only need $r$ to outgrow every fixed
power of $p$. Here $r=2^{\Theta(b^{1-\epsilon})}$ and $p=\Theta(b^{1+x})$, so that still holds.

\section{Gaussian row}

For each axis take $Q=p$. Then $\alpha^2/\theta\approx4d\alpha^2\le b'/2$, so $P=O(p)$.
Also $w^2=O(p/\alpha^2)=O(d)$, so $\theta w^3=O(d^{1/2})$. Summing over $d$ axes gives a
normalised cost of $O(d^{3/2}p^\delta)$. This is the exponent $\tfrac32\epsilon$, measured in
powers of $\log n$, used in the certificate.


\section{Sub-blocks with recentred chirps}

\begin{lemma}[Recentred blocks]\label{lem:recentred}
Let $\ell,j$ and a centre $c$ lie in one segment. Put $u=\ell-c$, $v=j-c$ and $h=j-\ell$. Then
\[
 X_{\ell,j}=\sigma h^2+2\beta_ch+\theta(v^2-u^2).
\]
Hence the principal block of $N$ on any run of consecutive indices inside a segment
starting at $c$ equals $D_cT_cD_c^{-1}$, with $D_c=\operatorname{diag}(e^{\pi\alpha^2\theta u^2})$ and
$T_c$ the Toeplitz matrix with symbol $e^{-\pi\alpha^2(\sigma h^2+2\beta_ch)}$.
\end{lemma}
\begin{proof}
In a segment $\beta_\ell=\beta_c+\theta u$ and $\beta_j=\beta_c+\theta v$. Expand
$X=(\sigma h+\beta_\ell)^2-\beta_j^2$, use $\sigma-\theta=1$ and $2hu=v^2-u^2-h^2$.
\end{proof}
This was checked in exact rational arithmetic for three centres per pair on all five test cases.

\paragraph{One Toeplitz inverse serves all blocks.} The symbol of $T_c$ is
$t_h\rho_c^{h}$, with $t_h=e^{-\pi\alpha^2(\sigma h^2-h)}$ and $\rho_c=e^{-\pi\alpha^2(2\beta_c+1)}$.
So $T_c=G_cT_\lambda G_c^{-1}$ with $G_c=\operatorname{diag}(\rho_c^{-i})$, and
$T_c^{-1}=G_cT_\lambda^{-1}G_c^{-1}$. The Gohberg--Semencul vectors of $T_c$ are the vectors
of $T_\lambda$, the block-length Toeplitz matrix with symbol $t_h$, scaled entrywise by
$\rho_c^{\mp i}$. They are computed once per line and rescaled in $O(\lambda)$ per block.
The scaled lower vectors decay like $e^{-\pi\alpha^2(\sigma+2\beta_c)i}$, and the upper ones like
$e^{-\pi\alpha^2(\sigma-2\beta_c)i}$. Since $|\beta_c|\le1/2$ and $\sigma-1=\theta>0$, every scaled
entry has modulus at most one, and $\|T_c-I\|<2e^{-\pi\alpha^2\theta}<0.1$.

\paragraph{Cuts.} Cut every segment into blocks of length at most $\lambda$, and treat each cut
like a wrap. Coupling entries across a cut have $n=h$, so $X\ge|h|(|h|-1)$. The same
$w\times w$ corner blocks therefore suffice. All block solves use principal blocks of $N=I+E$,
hence $\|M^{-1}\|\le2$. The Woodbury capacitance is cyclic block tridiagonal as before,
with one block per cut or wrap.

\paragraph{Precision and cost.} Inside a block, the chirp $D_c$ spans at most
$e^{\pi\alpha^2\theta\lambda^2}$, that is $4.53\,\alpha^2\theta\lambda^2$ bits. Choose
\[
 \lambda=\Bigl\lfloor\sqrt{Q/(4.53\,\alpha^2\theta)}\Bigr\rfloor,
\]
so that the precision reserve is at most $Q$. Couplings occur with density
$1/\lambda+\theta$ per output, each costing $O(w^3)$ with $w^2=O(Q/\alpha^2)$.

\paragraph{Reduced oversampling.} Take $\theta_i\approx\eta=1/(4dL^y)$ in prime selection
(Section~1 of the stack notes), so the padding factor is $\prod_i(1+\theta_i)=1+o(1)$.
Take $\alpha^2=\lfloor Q/(48d)\rfloor$ with $Q=\Theta(p)$. Then $\gamma=2d\alpha^2\le b'/4$, and
$\alpha^2\theta\ge1$ once $Q\ge192\,d^2L^y$, which holds when $2\epsilon+y<1+x$. This gives
$w^2=O(d)$ and $\lambda=\Theta(dL^{y/2})$, so the coupling costs
$O(w^3(1/\lambda+\theta))=O(d^{1/2}L^{-y/2}+d^{1/2}L^{-y})=O(1)$ per output for $y=\epsilon$.
Including the $O(1)$ chirped correlations, each axis costs $O(p^{1+\delta})$ per output. The
Gaussian row is therefore $O(d\,p^\delta)$: exponent $\epsilon+(1+x)\delta$.

\paragraph{Numerical check.} \texttt{scripts/check_subblock_inverse.py} cuts the
segments of the $(60,67)$, $\alpha^2=9$ instance into blocks of length $2$ and $3$, solves each
block by its own recentred Toeplitz system at $120$ digits, couples all cuts and wraps by
Woodbury, and matches the direct solve to $3\cdot10^{-64}$, against a target of $2^{-200}$.

\end{document}
